\textbf{递归对象}

设点列 \(P=(p_1,\dots,p_n)\) 已按第 \(1\) 维非降序排序，且
\[
p_i=(x_{i,1},\dots,x_{i,k}).
\]
定义加权 \(k\) 维偏序计数
\[
D_k(P;\alpha,\beta)
=
\sum_{\substack{1\le i<j\le n\\x_{i,1}\le x_{j,1},\dots,x_{i,k}\le x_{j,k}}}
\alpha_i\beta_j.
\]
原始点对计数只是特例 \(\alpha_i=\beta_i=1\)。

把 \(P\) 按下标分成左右两半 \(P_L,P_R\)。删去每个点的第一维，得到
\[
\widetilde P=(\pi(p_1),\dots,\pi(p_n)),\qquad \pi(p_i)=(x_{i,2},\dots,x_{i,k}).
\]
再定义新权值
\[
\alpha_i'=
\begin{cases}
\alpha_i,& p_i\in P_L,\\
0,& p_i\in P_R,
\end{cases}
\qquad
\beta_i'=
\begin{cases}
0,& p_i\in P_L,\\
\beta_i,& p_i\in P_R.
\end{cases}
\]

由于整体已按第一维排好序，任意跨区间点对自动满足第一维约束，因此有
\[
D_k(P;\alpha,\beta)
=
D_k(P_L;\alpha,\beta)
+
D_k(P_R;\alpha,\beta)
+
D_{k-1}(\widetilde P;\alpha',\beta').
\]

所以 CDQ 严格递归的不是“原始计数问题”，而是问题族 \(\{D_k\}\)。
对固定维度 \(k\) 看，它是“分治 + 降维”；对整个问题族看，它才是闭合递归。


\textbf{复杂度}

设 \(T_k(n)\) 为求解 \(D_k\) 在 \(n\) 个点上的最坏复杂度，则
\[
T_k(n)=2T_k(n/2)+T_{k-1}(n)+O(n),
\]
其中 \(O(n)\) 是当前层拆分、归并和辅助维护的线性代价。

一维基例为
\[
D_1(P;\alpha,\beta)=\sum_{1\le i<j\le n}\alpha_i\beta_j,
\]
线性扫描即可计算，故
\[
T_1(n)=\Theta(n).
\]

若已知
\[
T_{k-1}(n)=\Theta(n\log^{k-2}n),
\]
则
\[
T_k(n)=2T_k(n/2)+\Theta(n\log^{k-2}n).
\]
由扩展主定理
\[
T(n)=aT(n/b)+\Theta(n^{\log_b a}\log^m n)
\]
取
\[
a=b=2,\qquad m=k-2,
\]
可得
\[
T_k(n)=\Theta(n\log^{k-1}n).
\]

因此对任意 \(k\ge 1\),
\[
T_k(n)=\Theta(n\log^{k-1}n).
\]
